\ifdefined\gkver\else\def\gkver{student}\fi
\documentclass[\gkver]{gaokaozhenti}
\begin{document}

\chapter{2013年天津理综}
%% paperType: 理综物理部分

\begin{enumerate}
\item
%% number: 1
%% paperName: 2013年天津理综第 1 题 6 分
%% typeId: 1
%% type: 单选题
%% chapter: 原子和原子核
%% point: 天然放射性
%% method: 无
%% score: 6
%% degree: 850
%% duplicateId: 0
%% body:
下列说法正确的是\xzanswer{C}

\fourchoices[answer=C]
{原子核发生衰变时要遵守电荷守恒和质量守恒的规律}
{$\alpha$ 射线、$\beta$ 射线、$\gamma$ 射线都是高速运动的带电粒子流}
{氢原子从激发态向基态跃迁只能辐射特定频率的光子}
{发生光电效应时光电子的动能只与入射光的强度有关}

%% answer: C
%% memo:
\memoanswer{
本题原卷及材料中未提供详解，待补充。
}

\item
%% number: 2
%% paperName: 2013年天津理综第 2 题 6 分
%% typeId: 1
%% type: 单选题
%% chapter: 运动和力
%% point: 动量和冲量
%% method: 无
%% score: 6
%% degree: 650
%% duplicateId: 0
%% body:
我国女子短道速滑队在今年世锦赛上实现女子 $3000\Um$ 接力三连冠。观察发现，“接棒”的运动员甲提前站在“交棒”的运动员乙前面，并且开始向前滑行，待乙追上甲时，乙猛推甲一把，使甲获得更大的速度向前冲出。在乙推甲的过程中，忽略运动员与冰面间在水平方向上的相互作用，则\xzanswer{B}

\onepicture[w=4.5cm]{figs/02.png}

\fourchoices[answer=B]
{甲对乙的冲量一定等于乙对甲的冲量}
{甲、乙的动量变化一定大小相等方向相反}
{甲的动能增加量一定等于乙的动能减少量}
{甲对乙做多少负功，乙对甲就一定做多少正功}

%% answer: B
%% memo:
\memoanswer{
本题原卷及材料中未提供详解，待补充。
}

\item
%% number: 3
%% paperName: 2013年天津理综第 3 题 6 分
%% typeId: 1
%% type: 单选题
%% chapter: 电磁感应
%% point: 电磁感应能量分析
%% method: 无
%% score: 6
%% degree: 650
%% duplicateId: 0
%% body:
如图所示，纸面内有一矩形导体闭合线框 $abcd$，$ab$ 边长大于 $bc$ 边长，置于垂直纸面向里、边界为 $MN$ 的匀强磁场外，线框两次匀速地完全进入磁场，两次速度大小相同，方向均垂直于 $MN$。第一次 $ab$ 边平行 $MN$ 进入磁场，线框上产生的热量为 $Q_{1}$，通过线框导体横截面的电荷量为 $q_{1}$；第二次 $bc$ 边平行 $MN$ 进入磁场。线框上产生的热量为 $Q_{2}$，通过线框导体横截面的电荷量为 $q_{2}$，则\xzanswer{A}

\onepicture[w=4.5cm]{figs/03.png}

\fourchoices[answer=A]
{$Q_{1}>Q_{2}$，$q_{1}=q_{2}$}
{$Q_{1}>Q_{2}$，$q_{1}>q_{2}$}
{$Q_{1}=Q_{2}$，$q_{1}=q_{2}$}
{$Q_{1}=Q_{2}$，$q_{1}>q_{2}$}

%% answer: A
%% memo:
\memoanswer{
本题原卷及材料中未提供详解，待补充。
}

\item
%% number: 4
%% paperName: 2013年天津理综第 4 题 6 分
%% typeId: 1
%% type: 单选题
%% chapter: 交流电
%% point: 变压器
%% method: 无
%% score: 6
%% degree: 650
%% duplicateId: 0
%% body:
普通的交流电流表不能直接接在高压输电线路上测量电流，通常要通过电流互感器来连接，图中电流互感器 $ab$ 一侧线圈的匝数较少，工作时电流为 $I_{ab}$，$cd$ 一侧线圈的匝数较多，工作时电流为 $I_{cd}$，为了使电流表能正常工作，则\xzanswer{B}

\onepicture[w=6cm]{figs/04.png}

\fourchoices[answer=B]
{$ab$ 接 $MN$、$cd$ 接 $PQ$，$I_{ab}<I_{cd}$}
{$ab$ 接 $MN$、$cd$ 接 $PQ$，$I_{ab}>I_{cd}$}
{$ab$ 接 $PQ$、$cd$ 接 $MN$，$I_{ab}<I_{cd}$}
{$ab$ 接 $PQ$、$cd$ 接 $MN$，$I_{ab}>I_{cd}$}

%% answer: B
%% memo:
\memoanswer{
本题原卷及材料中未提供详解，待补充。
}

\item
%% number: 5
%% paperName: 2013年天津理综第 5 题 6 分
%% typeId: 1
%% type: 单选题
%% chapter: 力物体的平衡
%% point: 三力平衡
%% method: 无
%% score: 6
%% degree: 450
%% duplicateId: 0
%% body:
如图所示，小球用细绳系住，绳的另一端固定于 $O$ 点。现用水平力 $F$ 缓慢推动斜面体，小球在斜面上无摩擦地滑动，细绳始终处于直线状态，当小球升到接近斜面顶端时细绳接近水平，此过程中斜面对小球的支持力 $F_{N}$ 以及绳对小球的拉力 $F_{T}$ 的变化情况是\xzanswer{D}

\onepicture[w=6cm]{figs/05.png}

\fourchoices[answer=D]
{$F_{N}$ 保持不变，$F_{T}$ 不断增大}
{$F_{N}$ 不断增大，$F_{T}$ 不断减小}
{$F_{N}$ 保持不变，$F_{T}$ 先增大后减小}
{$F_{N}$ 不断增大，$F_{T}$ 先减小后增大}

%% answer: D
%% memo:
\memoanswer{
本题原卷及材料中未提供详解，待补充。
}

\item
%% number: 6
%% paperName: 2013年天津理综第 6 题 6 分
%% typeId: 2
%% type: 多选题
%% chapter: 电场
%% point: 带电粒子的往复运动
%% method: 无
%% score: 6
%% degree: 650
%% duplicateId: 0
%% body:
两个带等量正电的点电荷，固定在图中 $P$、$Q$ 两点，$MN$ 为 $PQ$ 连线的中垂线，交 $PQ$ 于 $O$ 点，$A$ 点为 $MN$ 上的一点。一带负电的试探电荷 $q$，从 $A$ 点由静止释放，只在静电力作用下运动。取无限远处的电势为零，则\xzanswer{BC}

\onepicture[w=4cm]{figs/06.png}

\fourchoices[answer=BC]
{$q$ 由 $A$ 向 $O$ 的运动是匀加速直线运动}
{$q$ 由 $A$ 向 $O$ 运动的过程电势能逐渐减小}
{$q$ 运动到 $O$ 点时的动能最大}
{$q$ 运动到 $O$ 点时电势能为零}

%% answer: BC
%% memo:
\memoanswer{
本题原卷及材料中未提供详解，待补充。
}

\item
%% number: 7
%% paperName: 2013年天津理综第 7 题 6 分
%% typeId: 2
%% type: 多选题
%% chapter: 机械振动
%% point: 振动图象
%% method: 无
%% score: 6
%% degree: 650
%% duplicateId: 0
%% body:
一列简谐横波沿直线传播，该直线上平衡位置相距 $9\Um$ 的 $a$、$b$ 两质点的振动图象如右图所示。下列描述该波的图象可能正确的是\xzanswer{AC}

\onepicture[align=right, w=4cm]{figs/07.png}

\fourchoices[answer=AC, ispicture=true, h=1.8cm]
{figs/07a.png}
{figs/07b.png}
{figs/07c.png}
{figs/07d.png}

%% answer: AC
%% memo:
\memoanswer{
本题原卷及材料中未提供详解，待补充。
}

\item
%% number: 8
%% paperName: 2013年天津理综第 8 题 6 分
%% typeId: 2
%% type: 多选题
%% chapter: 光学
%% point: 光的折射
%% method: 无
%% score: 6
%% degree: 450
%% duplicateId: 0
%% body:
固定的半圆形玻璃砖的横截面如图，$O$ 点为圆心，$OO'$ 为直径 $MN$ 的垂线。足够大的光屏 $PQ$ 紧靠玻璃砖右侧且垂直于 $MN$，由 $A$、$B$ 两种单色光组成的一束光沿半径方向射向 $O$ 点，入射光线与 $OO'$ 夹角 $\theta$ 较小时，光屏 $NQ$ 区域出现两个光斑，逐渐增大 $\theta$ 角，当 $\theta=\alpha$ 时，光屏 $NQ$ 区域 $A$ 光的光斑消失，继续增大 $\theta$ 角，当 $\theta=\beta$ 时，光屏 $NQ$ 区域 $B$ 光的光斑消失，则\xzanswer{AD}

\onepicture[w=3.5cm]{figs/08.png}

\fourchoices[answer=AD]
{玻璃砖对 $A$ 光的折射率比对 $B$ 光的大}
{$A$ 光在玻璃砖中传播速度比 $B$ 光的大}
{$\alpha<\theta<\beta$ 时，光屏上只有 1 个光斑}
{$\beta<\theta<\pi/2$ 时，光屏上只有 1 个光斑}

%% answer: AD
%% memo:
\memoanswer{
本题原卷及材料中未提供详解，待补充。
}

\item
%% number: 9
%% paperName: 2013年天津理综第 9 题 6 分
%% typeId: 4
%% type: 实验题
%% chapter: 电路
%% point: 伏安法测电阻
%% method: 无
%% score: 6
%% degree: 650
%% duplicateId: 0
%% body:
\begin{enumerate}
	\item
	“嫦娥一号”和“嫦娥二号”卫星相继完成了对月球的环月飞行，标志着我国探月工程的第一阶段已经完成。设“嫦娥二号”卫星环绕月球的运动为匀速圆周运动，它距月球表面的高度为 $h$，已知月球的质量为 $M$、半径为 $R$，引力常量为 $G$，则卫星绕月球运动的向心加速度 $a=$\tkanswer[2.2cm]{$\frac{GM}{(R+h)^{2}}$}，线速度 $v=$\tkanswer[2.2cm]{$\sqrt{\frac{GM}{R+h}}$}。

	\item
	某实验小组利用图示的装置探究加速度与力、质量的关系。

	\onepicture[align=right, w=6cm]{figs/09a.png}

	\begin{enumerate}
		\item
		下列做法正确的是\tkanswer{AD}（填字母代号）

		（A）调节滑轮的高度，使牵引木块的细绳与长木板保持平行

		（B）在调节木板倾斜度平衡木块受到的滑动摩擦力时，将装有砝码的砝码桶通过定滑轮拴木块上

		（C）实验时，先放开木块再接通打点计时器的电源

		（D）通过增减木块上的砝码改变质量时，不需要重新调节木板倾斜度
		\item
		为使砝码桶及桶内砝码的总重力在数值上近似等于木块运动时受到的拉力，应满足的条件是砝码桶及桶内砝码的总质量\tkanswer[1.4cm]{远小于}木块和木块上砝码的总质量（填远大于，远小于，或近似于）。
		\item
		甲、乙两同学在同一实验室，各取一套图示的装置放在水平桌面上，木块上均不放砝码，在没有平衡摩擦力的情况下，研究加速度 $a$ 与拉力 $F$ 的关系，分别得到图中甲、乙两条直线。设甲、乙用的木块质量分别为 $m_{\text{甲}}$、$m_{\text{乙}}$，甲、乙用的木块与木板间的动摩擦因数分别为 $\mu_{\text{甲}}$、$\mu_{\text{乙}}$，由图可知，$m_{\text{甲}}$\tkanswer[1cm]{小于}$m_{\text{乙}}$，$\mu_{\text{甲}}$\tkanswer[1cm]{大于}$\mu_{\text{乙}}$（填“大于”、“小于”或“等于”）。

		\onepicture[align=right, w=3cm]{figs/09b.png}
	\end{enumerate}
	\item
	要测绘一个标有“$3\UV$ $0.6\UW$”小灯泡的伏安特性曲线，灯泡两端的电压需要由零逐渐增加到 $3\UV$，并便于操作。已选用的器材有：

	电池组（电动势为 $4.5\UV$，内阻约 $1\UO$）；

	电流表（量程为 $0\sim250\UmA$，内阻约 $5\UO$）；

	电压表（量程为 $0\sim3\UV$，内阻约 $3\UkO$）；

	电键一个、导线若干。

	\begin{enumerate}
		\item
		实验中所用的滑动变阻器应选下列中的\tkanswer{A}（填字母代号）。

		（A）滑动变阻器（最大阻值 $20\UO$，额定电流 $1\UA$）

		（B）滑动变阻器（最大阻值 $1750\UO$，额定电流 $0.3\UA$）
		\item
		实验的电路图应选用下列的图\tkanswer{B}（填字母代号）。

		\fourchoices[answer=B, ispicture=true, h=1.8cm]
		{figs/09c.png}
		{figs/09d.png}
		{figs/09e.png}
		{figs/09f.png}
		\item
		实验得到小灯泡的伏安特性曲线如图所示。如果将这个小灯泡接到电动势为 $1.5\UV$，内阻为 $50\UO$ 的电源两端，小灯泡消耗的功率是\tkanswer[1cm]{0.1}$\UW$。

		\onepicture[align=right, w=4cm]{figs/09g.png}
	\end{enumerate}
\end{enumerate}

%% answer:
\jdanswer{
\begin{enumerate}
	\item
	$a=\frac{GM}{(R+h)^{2}}$，$v=\sqrt{\frac{GM}{R+h}}$
	\item
	① AD；② 远小于；③ 小于，大于
	\item
	① A；② B；③ $0.1\UW$
\end{enumerate}
}
%% memo:
\memoanswer{
本题原卷及材料中未提供详解，待补充。
}

\item
%% number: 10
%% paperName: 2013年天津理综第 10 题 16 分
%% typeId: 6
%% type: 计算题
%% chapter: 运动和力
%% point: 牛顿定律的应用
%% method: 无
%% score: 16
%% degree: 650
%% duplicateId: 0
%% body:
质量为 $m=4\Ukg$ 的小物块静止于水平地面上的 $A$ 点，现用 $F=10\UN$ 的水平恒力拉动物块一段时间后撤去，物块继续滑动一段位移停在 $B$ 点，$A$、$B$ 两点相距 $x=20\Um$，物块与地面间的动摩擦因数 $\mu=0.2$，$g$ 取 $10\Umsq$，求：
\begin{enumerate}
	\item
	物块在力 $F$ 作用过程发生位移 $x_{1}$ 的大小；
	\item
	撤去力 $F$ 后物块继续滑动的时间 $t$。
\end{enumerate}

%% answer:
\jdanswer{
\begin{enumerate}
	\item
	$16\Um$
	\item
	$2\Us$
\end{enumerate}
}
%% memo:
\memoanswer{
（1）设物块受到的滑动摩擦力为 $F_{f}$，则

$F_{f}=\mu mg$①

根据动能定理，对物块由 $A$ 到 $B$ 整个过程，有

$Fx_{1}-F_{f}x=0$②

代入数据，解得：$x_{1}=16\Um$③

（2）设撤去力 $F$ 时物块的速度为 $v$，此后物块的加速度为 $a$，滑动的位移为 $x_{2}$，则 $x_{2}=x-x_{1}$④

由牛顿第二定律得

$a=\frac{F_{f}}{m}$⑤

由匀变速直线运动公式得

$v^{2}=2ax_{2}$⑥

以物块运动的方向为正方向，由动量定理，得

$-F_{f}t=0-mv$⑦

代入数据，解得：$t=2\Us$⑧
}

\item
%% number: 11
%% paperName: 2013年天津理综第 11 题 18 分
%% typeId: 6
%% type: 计算题
%% chapter: 磁场
%% point: 带电粒子在磁场中的运动
%% method: 无
%% score: 18
%% degree: 450
%% duplicateId: 0
%% body:
一圆筒的横截面如图所示，其圆心为 $O$。筒内有垂直于纸面向里的匀强磁场，磁感应强度为 $B$。圆筒下面有相距为 $d$ 的平行金属板 $M$、$N$，其中 $M$ 板带正电荷，$N$ 板带等量负电荷。质量为 $m$、电荷量为 $q$ 的带正电粒子自 $M$ 板边缘的 $P$ 处由静止释放，经 $N$ 板的小孔 $S$ 以速度 $v$ 沿半径 $SO$ 方向射入磁场中，粒子与圆筒发生两次碰撞后仍从 $S$ 孔射出，设粒子与圆筒碰撞过程中没有动能损失，且电荷量保持不变，在不计重力的情况下，求：
\begin{enumerate}
	\item
	$M$、$N$ 间电场强度 $E$ 的大小；
	\item
	圆筒的半径 $R$；
	\item
	保持 $M$、$N$ 间电场强度 $E$ 不变，仅将 $M$ 板向上平移 $\frac{2}{3}d$，粒子仍从 $M$ 板边缘的 $P$ 处由静止释放，粒子自进入圆筒至从 $S$ 孔射出期间，与圆筒的碰撞次数 $n$。
\end{enumerate}

\onepicture[align=right, w=5cm]{figs/11.png}

%% answer:
\jdanswer{
\begin{enumerate}
	\item
	$E=\frac{mv^{2}}{2qd}$
	\item
	$R=\frac{\sqrt{3}mv}{3qB}$
	\item
	$n=3$
\end{enumerate}
}
%% memo:
\memoanswer{
（1）设两极板间的电压为 $U$，由动能定理得

$qU=\frac{1}{2}mv^{2}$①

由匀强电场中电势差与电场强度的关系得

$U=Ed$②

联立上式可得 $E=\frac{mv^{2}}{2qd}$③

（2）粒子进入磁场后做匀速圆周运动，运用几何关系做出圆心 $O'$，圆半径为 $r$，设第一次碰撞点为 $A$，由于粒子与圆筒发生两次碰撞又从 $S$ 孔射出，因此 $SA$ 弧所对圆心角 $\angle AO'S=\frac{\pi}{3}$。

由几何关系得 $r=R\tan\frac{\pi}{3}$④

粒子运动过程中洛伦兹力充当向心力，由牛顿第二定律，得

$qvB=m\frac{v^{2}}{r}$⑤

联立④⑤式得 $R=\frac{\sqrt{3}mv}{3qB}$⑥

\onepicture[w=5cm]{figs/11_an.jpg}

（3）保持 $M$、$N$ 间电场强度 $E$ 不变，$M$ 板向上平移 $\frac{2}{3}d$ 后，设板间电压为 $U'$，则 $U'=\frac{Ed}{3}=\frac{U}{3}$⑦

设粒子进入 $S$ 孔时的速度为 $v'$，由①式看出 $\frac{U'}{U}=\frac{v'^{2}}{v^{2}}$，结合⑦式可得

$v'=\frac{\sqrt{3}}{3}v$⑧

设粒子做圆周运动的半径为 $r'$，则

$r'=\frac{\sqrt{3}mv}{3qB}$⑨

设粒子从 $S$ 到第一次与圆筒碰撞期间的轨道所对圆心角为 $\theta$，比较⑥⑨两式得到 $r'=R$，可见 $\theta=\frac{\pi}{2}$⑩

粒子须经过这样的圆弧才能从 $S$ 孔射出，故

$n=3$\Circled{11}
}

\item
%% number: 12
%% paperName: 2013年天津理综第 12 题 20 分
%% typeId: 6
%% type: 计算题
%% chapter: 功和能
%% point: 动能定理
%% method: 无
%% score: 20
%% degree: 450
%% duplicateId: 0
%% body:
超导现象是20世纪人类重大发现之一，日前我国已研制出世界传输电流最大的高温超导电缆并成功示范运行。
\begin{enumerate}
	\item
	超导体在温度特别低时电阻可以降到几乎为零，这种性质可以通过实验研究。将一个闭合超导金属圈环水平放置在匀强磁场中，磁感线垂直于圈环平面向上，逐渐降低温度使环发生由正常态到超导态的转变后突然撤去磁场，若此后环中的电流不随时间变化，则表明其电阻为零。请指出自上往下看环中电流方向，并说明理由。
	\item
	为探究该圆环在超导状态的电阻率上限 $\rho$，研究人员测得撤去磁场后环中电流为 $I$，并经一年以上的时间 $t$ 未检测出电流变化。实际上仪器只能检测出大于 $\Delta I$ 的电流变化，其中 $\Delta I\propto I$，当电流的变化小于 $\Delta I$ 时，仪器检测不出电流的变化，研究人员便认为电流没有变化。设环的横截面积为 $S$，环中定向移动电子的平均速率为 $v$，电子质量为 $m$、电荷量为 $e$。试用上述给出的各物理量，推导出 $\rho$ 的表达式。
	\item
	若仍使用上述测量仪器，实验持续时间依旧为 $t$，为使实验获得的该圆环在超导状态的电阻率上限 $\rho$ 的准确程度更高，请提出你的建议，并简要说明实现方法。
\end{enumerate}

%% answer:
\jdanswer{
\begin{enumerate}
	\item
	逆时针方向。撤去磁场瞬间，环所围面积的磁通量突变为零，由楞次定律可知，环中电流的磁场方向应与原磁场方向相同，即向上。由右手螺旋定则可知，环中电流的方向是沿逆时针方向。
	\item
	$\rho=\frac{mvS\Delta I}{etI^{2}}$
	\item
	由 $\rho=\frac{mvS\Delta I}{etI^{2}}$ 看出，在题设条件限制下，适当增大超导电流，可以使实验获得 $\rho$ 的准确程度更高，通过增大穿过该环的磁通量变化率可实现增大超导电流。
\end{enumerate}
}
%% memo:
\memoanswer{
（1）逆时针方向。撤去磁场瞬间，环所围面积的磁通量突变为零，由楞次定律可知，环中电流的磁场方向应与原磁场方向相同，即向上。由右手螺旋定则可知，环中电流的方向是沿逆时针方向。

（2）设圆环周长为 $l$、电阻为 $R$，由电阻定律得

$R=\rho\frac{l}{S}$①

设 $t$ 时间内环中电流释放焦耳热而损失的能量为 $\Delta E$，由焦耳定律得

$\Delta E=I^{2}Rt$②

设环中单位体积内定向移动电子数为 $n$，则

$I=nevS$③

式中 $n$、$e$、$S$ 不变，只有定向移动电子的平均速率的变化才会引起环中电流的变化，电流变化大小取 $\Delta I$ 时，相应定向移动电子的平均速率的变化大小为 $\Delta v$，则

$\Delta I=neS\Delta v$④

设环中定向移动电子减少的动能总和为 $\Delta E_{k}$，则

$\Delta E_{k}=nlS\left[\frac{1}{2}mv^{2}-\frac{1}{2}m(v-\Delta v)^{2}\right]$⑤

由于 $\Delta I\ll I$，可得

$\Delta E_{k}=\frac{lmv}{e}\Delta I$⑥

根据能量守恒定律，得

$\Delta E=\Delta E_{k}$⑦

联立上述各式，得 $\rho=\frac{mvS\Delta I}{etI^{2}}$⑧

（3）由 $\rho=\frac{mvS\Delta I}{etI^{2}}$ 看出，在题设条件限制下，适当增大超导电流，可以使实验获得 $\rho$ 的准确程度更高，通过增大穿过该环的磁通量变化率可实现增大超导电流。
}

\end{enumerate}

\end{document}
